\documentclass[11pt]{article}
\usepackage[margin=1in]{geometry}
\usepackage{booktabs}
\usepackage{graphicx}
\usepackage{float}
\usepackage{xcolor}
\usepackage{hyperref}
\usepackage{enumitem}
\usepackage{tabularx}
\usepackage{multirow}
\usepackage{caption}

\definecolor{darkgreen}{rgb}{0.0, 0.5, 0.0}
\definecolor{darkred}{rgb}{0.7, 0.0, 0.0}
\newcommand{\done}[1]{\textcolor{darkgreen}{\textbf{#1}}}
\newcommand{\pending}[1]{\textcolor{darkred}{\textbf{#1}}}
\newcommand{\infs}{inf/s}

\title{MiniRocketHLS Resubmission Summary:\\
\large What Changed Since FCCM 2026 Rejection}
\author{Prepared for Philip --- Data \& Analysis Summary}
\date{April 3, 2026}

\begin{document}
\maketitle

%% ============================================================
\section{Executive Summary}
%% ============================================================

FCCM 2026 paper \#238 (scores: 2, 3, 3, 4) was rejected on five grounds:
no GPU baseline, Python-only CPU comparison, limited scalability analysis,
questioned novelty, and missing power data.
\textbf{All five complaints have been addressed with new experimental data.}

Since rejection, we have completed:
\begin{itemize}[nosep]
  \item Tesla T4 GPU baselines for both MiniRocket and HYDRA (batch=256 \emph{and} batch=1)
  \item C++ \texttt{-O3} single-threaded CPU baselines for MiniRocket (3 datasets)
  \item 8-point scalability sweep (length 64--8192) for both GPU and FPGA
  \item CU scaling analysis (1/2/3 compute units, 4 datasets)
  \item Measured power: FPGA 26\,W (xbutil), GPU 46--70\,W (nvidia-smi), CPU 90\,W (TDP)
  \item Fused CONV+PPV kernel achieving 0 DSP and 68\% BRAM reduction
  \item HYDRA \texttt{ap\_fixed<32,16>} kernel validated on hardware (45--70$\times$ over CPU)
\end{itemize}

The paper is now a \textbf{three-algorithm, two-platform comparison} with honest
GPU numbers, scalability analysis, and design insight --- substantially stronger
than the rejected submission.

%% ============================================================
\section{Reviewer Complaints $\rightarrow$ New Evidence}
%% ============================================================

\begin{table}[H]
\centering
\caption{Mapping reviewer complaints to new experimental evidence.}
\label{tab:complaints}
\small
\begin{tabularx}{\textwidth}{p{0.22\textwidth} p{0.38\textwidth} p{0.32\textwidth}}
\toprule
\textbf{Complaint} & \textbf{New Evidence} & \textbf{Key Numbers} \\
\midrule

\textbf{1. No GPU baseline} &
Tesla T4 comparison for MiniRocket (3 datasets) and HYDRA (3 datasets).
Both batch=256 and batch=1 measured.
Power recorded via \texttt{nvidia-smi}. &
GPU wins batched: 1.5--6.1$\times$ (MiniRocket), 2.2--2.8$\times$ (HYDRA).
FPGA wins batch=1: 4.2--7.5$\times$ lower latency. \\
\addlinespace

\textbf{2. Python-only CPU} &
C++ \texttt{-O3} single-threaded baselines for MiniRocket on all 3 benchmark datasets.
Per-sample latency with P50/P95/P99 statistics. &
C++ throughput: 48--289~\infs{}.
FPGA speedup over C++: 8.5--58.7$\times$ (3\,CU). \\
\addlinespace

\textbf{3. Limited scalability} &
8-point sweep (64--8192) for both GPU and FPGA.
CU scaling: 1/2/3 CU on 4 datasets.
GPU advantage shrinks dramatically with series length. &
GPU/FPGA ratio drops 65$\times \rightarrow$ 3.8$\times$ as length grows 64$\rightarrow$8192.
At batch=1, FPGA matches GPU at all lengths. \\
\addlinespace

\textbf{4. Novelty questioned} &
Fused CONV+PPV kernel eliminates intermediate array, achieves 0~DSP via ternary weight exploitation.
Three algorithms (MiniRocket + HYDRA + MultiRocket) on one FPGA platform. &
0~DSP (vs.\ 724 separate), BRAM $-68\%$.
Throughput gain: 6.6--10.9$\times$ from fusion alone. \\
\addlinespace

\textbf{5. Missing power data} &
Measured: FPGA 26\,W (\texttt{xbutil}), GPU 46--70\,W (\texttt{nvidia-smi}), CPU 90\,W (TDP).
Energy-per-inference computed for all configurations. &
FPGA energy: 5.2--35.0~mJ/inf.
GPU energy: 1.8--21.1~mJ/inf.
FPGA wins at long series: 9.2 vs.\ 15.7~mJ/inf (MosquitoSound). \\

\bottomrule
\end{tabularx}
\end{table}

\paragraph{Honesty note.}
We do not claim FPGA beats GPU on batched throughput --- it does not.
At batch=256, GPU T4 is 1.5--6.1$\times$ faster for MiniRocket.
Both platforms are PCIe-attached and host-dispatched; neither operates in a true streaming mode in this evaluation.
The FPGA advantage at batch=1 is real but partly reflects PyTorch dispatch overhead on the GPU side.
We lead with \emph{design insights} and \emph{deployment tradeoffs}, not raw speedup claims.

%% ============================================================
\section{Key Results}
%% ============================================================

The following tables contain hardware-validated data from the Alveo U280 and
Tesla T4 GPU. All FPGA numbers are from actual \texttt{.xclbin} execution,
not simulation.

\subsection{GPU vs.\ FPGA Throughput (MiniRocket)}

\begin{table}[H]
\centering
\caption{GPU (Tesla T4) vs.\ FPGA (Alveo U280) throughput and energy comparison
for MiniRocket feature extraction. FPGA design uses $<$5\% LUT, 0 DSP.}
\label{tab:gpu_comparison}
\begin{tabular}{lrrrrr}
\toprule
\textbf{Dataset} & \textbf{Length} & \textbf{GPU T4} & \textbf{FPGA 3CU} & \textbf{Ratio} & \textbf{Energy Ratio} \\
 & & (inf/s) & (inf/s) & (GPU/FPGA) & (GPU/FPGA) \\
\midrule
InsectSound   & 600  & 30{,}469 & 5{,}012 & 6.1$\times$ & 2.9$\times$ \\
MosquitoSound & 3750 & 4{,}198  & 2{,}820 & 1.5$\times$ & 1.7$\times$ \\
FruitFlies    & 5000 & 3{,}227  & 2{,}191 & 1.5$\times$ & 1.8$\times$ \\
\midrule
\multicolumn{2}{l}{\textbf{Power (under load)}} & \textbf{54--68\,W} & \textbf{26\,W} & & \\
\bottomrule
\end{tabular}

\vspace{0.5em}

\begin{minipage}{\textwidth}
\footnotesize
Energy ratio = (GPU mJ/inf) / (FPGA mJ/inf). Values $>$1 favor FPGA.
FPGA power is static-dominated ($<$5\% utilization); GPU power scales with load.
\end{minipage}
\end{table}

\subsection{Single-Sample Latency}

\begin{table}[H]
\centering
\caption{Single-sample inference latency: GPU vs.\ FPGA. Both platforms use
PCIe-attached host dispatch. The FPGA's lower per-invocation overhead (PCIe
DMA vs.\ CUDA runtime) yields 4--8$\times$ lower latency for MiniRocket.
Optimized GPU implementations (fused CUDA kernels, CUDA graphs) could
reduce this gap; results reflect standard PyTorch + cuDNN.}
\label{tab:batch1_latency}
\begin{tabular}{l r r r r r}
\toprule
\textbf{Dataset} & \textbf{Length} & \textbf{GPU P50} & \textbf{FPGA 3CU} & \textbf{Latency} & \textbf{Throughput} \\
 & & (ms) & (ms) & \textbf{Ratio} & \textbf{Ratio} \\
\midrule
\multicolumn{6}{l}{\textit{Single-sample (batch\,=\,1)}} \\
InsectSound   & 600   & 1.503 & 0.200 & 7.5$\times$ & 8.2$\times$ \\
MosquitoSound & 3{,}750 & 2.374 & 0.355 & 6.7$\times$ & 6.9$\times$ \\
FruitFlies    & 5{,}000 & 1.911 & 0.456 & 4.2$\times$ & 4.4$\times$ \\
\midrule
\multicolumn{6}{l}{\textit{Batched (batch\,=\,256)}} \\
InsectSound   & 600   & 0.037 & 0.200 & 0.18$\times$ & 0.18$\times$ \\
MosquitoSound & 3{,}750 & 0.233 & 0.355 & 0.66$\times$ & 0.66$\times$ \\
FruitFlies    & 5{,}000 & 0.306 & 0.456 & 0.67$\times$ & 0.67$\times$ \\
\bottomrule
\end{tabular}

\vspace{0.3em}

\begin{minipage}{\textwidth}
\footnotesize
GPU P50 = median over 200 inferences with \texttt{torch.cuda.synchronize()}.
FPGA latency includes PCIe round-trip and all 3 pipeline stages (K1+K2+K3).
Latency ratio $>$1 favors FPGA.
\textbf{Note:} The FPGA always processes one sample per kernel invocation ---
there is no batched dispatch mode. The FPGA column is identical in both
sections; what changes is the GPU operating point. At batch=256, the GPU
amortizes CUDA launch overhead across 256 samples, reducing per-sample
latency below the FPGA's fixed per-invocation cost. At batch=1, both
platforms pay per-sample dispatch overhead, and the FPGA's lighter
PCIe DMA path wins.
\end{minipage}
\end{table}

\subsection{Scalability: GPU Advantage Shrinks with Series Length}

\begin{table}[H]
\centering
\caption{MiniRocket throughput scaling with time series length.
GPU advantage (batch\,=\,256) shrinks from 65$\times$ to 3.8$\times$ as
length increases from 64 to 8{,}192. At batch\,=\,1, the FPGA matches
or exceeds GPU throughput at all lengths. FPGA length is a runtime
parameter --- no rebuild required.}
\label{tab:scalability}
\resizebox{\textwidth}{!}{
\begin{tabular}{r r r r r r r}
\toprule
\textbf{Length} & \textbf{GPU b256} & \textbf{GPU b1} & \textbf{FPGA 1CU} & \textbf{GPU/FPGA} & \textbf{b1/FPGA} & \textbf{Dilations} \\
 & (inf/s) & (inf/s) & (inf/s) & (b256) & & \\
\midrule
64      & 174{,}912 & 1{,}204 & 2{,}681 & 65$\times$   & 0.45$\times$ & 3 \\
128     & 115{,}308 & 1{,}091 & 2{,}458 & 47$\times$   & 0.44$\times$ & 4 \\
256     & 65{,}826  & 950     & 2{,}512 & 26$\times$   & 0.38$\times$ & 5 \\
512     & 33{,}826  & 798     & 2{,}204 & 15$\times$   & 0.36$\times$ & 6 \\
1{,}024 & 20{,}178  & 692     & 1{,}745 & 12$\times$   & 0.40$\times$ & 7 \\
2{,}048 & 9{,}331   & 661     & 1{,}303 & 7.2$\times$  & 0.51$\times$ & 8 \\
4{,}096 & 4{,}249   & 606     & 794     & 5.4$\times$  & 0.76$\times$ & 9 \\
8{,}192 & 1{,}919   & 539     & 511     & 3.8$\times$  & 1.05$\times$ & 10 \\
\bottomrule
\end{tabular}}

\vspace{0.3em}

\begin{minipage}{\textwidth}
\footnotesize
GPU/FPGA $>$1 favors GPU; b1/FPGA $<$1 favors FPGA.
FPGA throughput scales as $O(1/n)$ with length (fused CONV\_LOOP at II\,=\,1).
At batch=1, GPU is launch-overhead dominated ($\sim$1\,ms fixed cost)
while FPGA latency is purely compute-proportional.
\end{minipage}
\end{table}

\subsection{HYDRA: GPU Advantage Much Smaller}

\begin{table}[H]
\centering
\caption{HYDRA inference: GPU (Tesla T4) vs.\ FPGA (Alveo U280, 1 CU).
HYDRA's 512 learned kernels with arbitrary float weights reduce the GPU's
batched-convolution advantage compared to MiniRocket (2--3$\times$ vs.\
6$\times$ at length 600). Both platforms use PCIe host dispatch.}
\label{tab:hydra_gpu}
\begin{tabular}{l r r r r r}
\toprule
\textbf{Dataset} & \textbf{Length} & \textbf{GPU b256} & \textbf{GPU b1} & \textbf{FPGA 1CU} & \textbf{FPGA/GPU} \\
 & & (inf/s) & (inf/s) & (inf/s) & (b256 / b1) \\
\midrule
InsectSound   & 600   & 13{,}895 & 56  & 6{,}326 & 0.46$\times$ / 113$\times$ \\
MosquitoSound & 3{,}750 & 5{,}139  & 53  & 1{,}937 & 0.38$\times$ / 36$\times$ \\
FruitFlies    & 5{,}000 & 4{,}267  & 57  & 1{,}507 & 0.35$\times$ / 26$\times$ \\
\midrule
\multicolumn{2}{l}{\textbf{Power (under load)}} & \textbf{54--70\,W} & & \textbf{26\,W} & \\
\bottomrule
\end{tabular}

\vspace{0.3em}

\begin{minipage}{\textwidth}
\footnotesize
The large batch=1 gap (26--113$\times$) reflects PyTorch dispatch overhead;
a fused CUDA kernel would reduce this substantially.
The batched comparison (0.35--0.46$\times$) is more representative of sustained throughput.
Energy at batch=256: GPU 3.9--16.4\,mJ/inf vs FPGA 4.1--17.2\,mJ/inf (near parity).
\end{minipage}
\end{table}

\subsection{CU Scaling Efficiency}

\begin{table}[H]
\centering
\caption{MiniRocket CU scaling with fused CONV+PPV kernel. 2-CU runs at 300 MHz (timing met), 3-CU at 269 MHz (timing failed, auto-throttled). Scaling efficiency ranges from 75--88\% for 2 CU and 59--94\% for 3 CU.}
\label{tab:cu_scaling}
\begin{tabular}{l r r r r r}
\toprule
Dataset & 1 CU & 2 CU & 3 CU & 2/1 Scaling & 3/1 Scaling \\
\midrule
GunPoint & 3{,}797 & 5{,}999 & 9{,}585 & 1.58$\times$ & 2.52$\times$ \\
InsectSound & 1{,}967 & 3{,}191 & 5{,}012 & 1.62$\times$ & 2.55$\times$ \\
MosquitoSound & 937 & 1{,}497 & 2{,}820 & 1.60$\times$ & 3.01$\times$ \\
FruitFlies & 742 & 1{,}121 & 2{,}191 & 1.51$\times$ & 2.95$\times$ \\
\bottomrule
\end{tabular}
\end{table}

\subsection{Fused Kernel Impact}

\begin{table}[H]
\centering
\caption{Impact of fused CONV+PPV kernel on MiniRocket (1 CU, 300 MHz). Fusion eliminates the intermediate \texttt{convolutions[16][8192]} array, achieving 0 DSP by exploiting ternary weights, and uses dynamic loop bounds that scale with time series length.}
\label{tab:fused}
\begin{tabular}{l r r r}
\toprule
& Separate & Fused & Improvement \\
\midrule
\multicolumn{4}{l}{\textit{Throughput (inferences/sec)}} \\
\quad GunPoint & 508 & 3{,}797 & 7.5$\times$ \\
\quad InsectSound & 296 & 1{,}967 & 6.6$\times$ \\
\quad MosquitoSound & 88 & 937 & 10.6$\times$ \\
\quad FruitFlies & 68 & 742 & 10.9$\times$ \\
\addlinespace
\multicolumn{4}{l}{\textit{Resource Utilization (HLS synthesis)}} \\
\quad LUTs & 48{,}812 & 57{,}090 & +17\% \\
\quad FFs & 88{,}492 & 94{,}283 & +7\% \\
\quad BRAM (18K) & 576 & 185 & -68\% \\
\quad DSP48E2 & 724 & 0 & -100\% \\
\addlinespace
\quad Intermediate array & $16 \times 8192$ floats & Eliminated & $-$100\% \\
\bottomrule
\end{tabular}
\end{table}

\subsection{Energy Efficiency}

\begin{table}[H]
\centering
\caption{Energy Efficiency: MiniRocket Fused Kernel --- CPU vs.\ GPU vs.\ FPGA}
\label{tab:power}
\small
\begin{tabular}{llrrrr}
\toprule
\textbf{Dataset} & \textbf{Platform} & \textbf{Throughput} & \textbf{Power} & \textbf{Energy} & \textbf{Eff.} \\
 & & (inf/s) & (W) & (mJ/inf) & \textbf{Ratio} \\
\midrule
InsectSound & CPU C++ (-O3) & 289 & 90\textsuperscript{a} & 311.4 & 1$\times$ \\
(length 600) & GPU T4 (PyTorch) & 30{,}469 & 54\textsuperscript{b} & 1.8 & 176$\times$ \\
 & FPGA 1-CU & 1{,}967 & 26\textsuperscript{c} & 13.2 & 24$\times$ \\
 & FPGA 3-CU & 5{,}012 & 26 & 5.2 & 60$\times$ \\
\midrule
MosquitoSound & CPU C++ (-O3) & 48 & 90 & 1{,}875.0 & 1$\times$ \\
(length 3750) & GPU T4 (PyTorch) & 4{,}198 & 66 & 15.7 & 119$\times$ \\
 & FPGA 1-CU & 937 & 26 & 27.7 & 68$\times$ \\
 & FPGA 3-CU & 2{,}820 & 26 & 9.2 & 204$\times$ \\
\midrule
FruitFlies & CPU C++ (-O3) & 258 & 90 & 348.8 & 1$\times$ \\
(length 5000) & GPU T4 (PyTorch) & 3{,}227 & 68 & 21.1 & 17$\times$ \\
 & FPGA 1-CU & 742 & 26 & 35.0 & 10$\times$ \\
 & FPGA 3-CU & 2{,}191 & 26 & 11.9 & 29$\times$ \\
\bottomrule
\end{tabular}

\vspace{0.3em}

\begin{minipage}{\textwidth}
\footnotesize
\textsuperscript{a}TDP upper bound (single socket); actual load power likely lower.
Measured CPU power requires reading Intel RAPL counters via
\texttt{/sys/class/powercap/intel-rapl/}, which needs root privileges
(\texttt{sudo}) that are unavailable on our shared server. If sudo access
can be arranged before submission, we can replace TDP with measured
load power.
\textsuperscript{b}Measured via \texttt{nvidia-smi} during inference.
\textsuperscript{c}Measured via \texttt{xbutil examine}; static-dominated at $<$5\% utilization.
\end{minipage}
\end{table}


%% ============================================================
\section{Revised Paper Narrative}
%% ============================================================

The paper should make four core arguments:

\begin{enumerate}[leftmargin=*]

\item \textbf{Algorithm--hardware co-design insight.}
MiniRocket's fixed ternary weights $\{-1, +2\}$ eliminate all DSP blocks
when exploited at the HLS level. Fusing convolution and pooling into a single
pass removes the intermediate \texttt{convolutions[16][8192]} array, cutting
BRAM by 68\% and delivering 6.6--10.9$\times$ throughput improvement.
The entire MiniRocket fused kernel uses \textbf{0~DSP} and \textbf{$<$5\% LUT}
on a U280 --- the FPGA is $>$95\% unutilized.
This is not a brute-force acceleration; it is a case study in matching
algorithm structure to hardware primitives.

\item \textbf{Batch-size--dependent tradeoff.}
GPU wins at batch$\geq$32 (amortized kernel launch, SIMT parallelism).
FPGA wins at batch=1 with \textbf{4.2--7.5$\times$ lower latency} because
there is no framework dispatch overhead and the pipeline is always warm.
This is the honest picture: the ``right'' accelerator depends on the
deployment scenario (batched analytics vs.\ real-time per-sample inference).

\item \textbf{Scalability favors FPGA as series grow.}
GPU's batched advantage shrinks from 65$\times$ (length~64) to 3.8$\times$ (length~8192).
At batch=1, FPGA matches or exceeds GPU at \emph{every} length tested.
For long time series (e.g., audio at 16\,kHz), the FPGA becomes increasingly competitive
even in batched mode, while consuming 2--3$\times$ less power.

\item \textbf{Three algorithms, one platform.}
MiniRocket, HYDRA, and MultiRocket share the Alveo U280 with different
kernel architectures (ternary fixed-weight vs.\ learned-weight convolution,
PPV vs.\ max+mean pooling). HYDRA's GPU advantage is only 2.2--2.8$\times$
(vs.\ 1.5--6.1$\times$ for MiniRocket) because 512 arbitrary-weight kernels
stress GPU caches differently. This multi-algorithm comparison on a single
platform is novel in the FPGA time series literature.

\end{enumerate}

%% ============================================================
\section{What We Can Still Add Before Submission}
%% ============================================================

\begin{table}[H]
\centering
\caption{Potential additions ranked by impact and effort.}
\label{tab:additions}
\small
\footnotesize
\begin{tabularx}{\textwidth}{Xlll}
\toprule
\textbf{Item} & \textbf{Impact} & \textbf{Effort} & \textbf{Status} \\
\midrule
MultiRocket full benchmarks (3 datasets) & High & 1--2 weeks & In progress \\
MultiRocket GPU baseline & High & 3 hours & Blocked \\
Roofline model (analytical) & Medium & 2--3 hours & Not started \\
Real UCR data on GPU (not synthetic) & Medium & 1--2 hours & Not started \\
Measured CPU power (RAPL) & Low & Needs sudo & Blocked \\
Literature survey of FPGA TSC papers & Medium & 4--6 hours & Not started \\
\bottomrule
\end{tabularx}
\end{table}

\paragraph{Recommendation.}
\textbf{Priority 1} (before submission): Complete MultiRocket benchmarks and GPU baselines
to deliver a true three-algorithm story. Without MultiRocket hardware numbers on all
3~benchmark datasets, the ``three algorithms on one platform'' claim is weakened.
\textbf{Priority 2}: Add a roofline model (purely analytical, no experiments) to
quantify how far each kernel is from the compute/memory bound.
\textbf{Priority 3}: Run GPU baselines on real UCR test sets instead of synthetic data,
eliminating a potential reviewer concern about data representativeness.

%% ============================================================
\section{Recommended Venue and Timeline}
%% ============================================================

\begin{table}[H]
\centering
\caption{Venue options with open deadlines (as of April 3, 2026).}
\label{tab:venues}
\small
\begin{tabular}{lllcl}
\toprule
\textbf{Venue} & \textbf{Deadline} & \textbf{Notification} & \textbf{Pages} & \textbf{Fit} \\
\midrule
\done{FPT 2026 Journal Track (TRETS)} & \textbf{May 1, 2026} & Aug 14, 2026 & 22 & Excellent \\
FPT 2026 Conference Track & Jun 12, 2026 & Aug 14, 2026 & $\sim$8 & Excellent \\
DATE 2027 & $\sim$Sep 2026 & $\sim$Nov 2026 & 6 & Good \\
ICCAD 2026 & Apr 14, 2026 & Jul 11, 2026 & 8 & Moderate \\
TRETS (rolling) & Anytime & 4--8 months & 22 & Excellent \\
ReConFig 2026 & $\sim$Aug 2026 & TBD & TBD & Good \\
FCCM 2027 & $\sim$Jan 2027 & $\sim$Mar 2027 & 8 & Excellent \\
\bottomrule
\end{tabular}
\end{table}

\subsection{Primary Recommendation: FPT 2026 Journal Track via TRETS}

\textbf{Deadline: May 1, 2026} (28 days from today). Submit to ACM TRETS under
the FPT~2026 Journal Track. Accepted papers are published as full journal articles
in TRETS (the premier FPGA journal) and presented at FPT~2026 in Honolulu
(December 7--10, 2026). The 22-page limit accommodates all three algorithms,
the full scalability sweep, GPU comparison, and design methodology discussion
without space pressure.

\textbf{Why FPT/TRETS:} (1)~FPGA-focused community that values design methodology
and honest evaluation over raw speedup claims; (2)~journal publication carries
more weight than conference proceedings; (3)~22~pages lets us tell the complete
story; (4)~presentation at a prestigious venue.

\subsection{Fallback: FPT 2026 Conference Track}

If May~1 is too tight, the regular conference track (deadline June~12) gives
6 more weeks. The tradeoff: $\sim$8-page conference paper in IEEE Xplore
instead of a TRETS journal article. Still an excellent venue.

\subsection{Long-term: DATE 2027 or FCCM 2027}

If FPT does not work out, DATE~2027 ($\sim$September deadline, 6~pages)
and FCCM~2027 ($\sim$January 2027, 8~pages) are strong options.
FCCM is the most natural venue for this work but requires waiting
until January.

\subsection{Suggested Timeline}

\begin{table}[H]
\centering
\caption{Suggested timeline for FPT Journal Track submission (May 1, 2026).}
\label{tab:timeline}
\small
\begin{tabular}{@{}l p{0.72\textwidth}@{}}
\toprule
\textbf{Date} & \textbf{Milestone} \\
\midrule
Apr 3--7 & Complete MultiRocket model retraining; run GPU baselines on real data \\
Apr 7--10 & Finalize all LaTeX tables; run MultiRocket FPGA benchmarks \\
Apr 10--14 & Philip drafts/revises paper narrative; Rishi provides roofline model \\
Apr 14--21 & Internal review; address gaps; polish figures \\
Apr 21--28 & Final revisions; co-author review \\
Apr 29--May 1 & Format check; submit to TRETS/FPT \\
\bottomrule
\end{tabular}
\end{table}

\end{document}
